% 2026, Prof. Dr. Matthias Jung (DL9MJ)
% FELDWELLENWIDERSTAND
\begin{tikzpicture}

\begin{axis}[
    name=mainaxis,
    width=\linewidth,
    height=0.66*\linewidth,
    xmode=log,
    ymode=log,
    xmin=0.01,
    xmax=100,
    ymin=1,
    ymax=100000,
    domain=0.01:100,
    samples=500,
    clip=false,
    xlabel={
        Distanz, normiert:
        $\displaystyle
        x=\frac{d}{\lambda/(2\pi)}
        $
    },
    ylabel={
        Feldwellenwiderstand:
        $\displaystyle
        Z_\mathrm{F}
        =
        \left|\frac{\underline E}{\underline H}\right|
        \text{ in }\unit{\ohm}
        $
    },
    tick label style={
        font=\normalsize
    },
    xtick={0.01,0.1,1,10,100},
    xticklabels={
        {$0{,}01$},
        {$0{,}1$},
        {$1$},
        {$10$},
        {$100$}
    },
    grid=both,
]

% Feldwellenwiderstand des elektrischen Dipols:
%
% Z_el = Z0 * sqrt[
%   (x^4 - x^2 + 1) /
%   (x^2 * (1 + x^2))
% ]
%
\addplot[
    DARCorange,
    very thick
]
{
    376.730313668
    *
    sqrt(
        (x^4 - x^2 + 1)
        /
        (x^2 * (1 + x^2))
    )
};

% Feldwellenwiderstand der magnetischen Loop:
%
% Z_mag = Z0 * sqrt[
%   (x^2 * (1 + x^2)) /
%   (x^4 - x^2 + 1)
% ]
%
\addplot[
    DARCgreen,
    very thick
]
{
    376.730313668
    *
    sqrt(
        (x^2 * (1 + x^2))
        /
        (x^4 - x^2 + 1)
    )
};

% Feldwellenwiderstand des freien Raumes
\draw[
    black,
    dashed,
    line width=0.9pt
]
(axis cs:0.01,376.730313668)
--
(axis cs:100,376.730313668);

\node[
    anchor=south west,
]
at (axis cs:0.02,400)
{
    $\displaystyle Z_0\approx\qty{377}{\ohm}$
};

% Grenze reaktives Nahfeld:
% r = lambda/(2 pi), also x = 1
\draw[
    black,
    dashed,
    line width=1pt
]
(axis cs:1,1)
--
(axis cs:1,100000);

\node[
    anchor=north west
]
at (axis cs:1.05,50000)
{
    $\displaystyle d=\frac{\lambda}{2\pi}$
};

% Grenze Fernfeld:
% r = 4 lambda entspricht x = 8 pi
\draw[
    black,
    dashed,
    line width=1pt
]
(axis cs:{8*pi},1)
--
(axis cs:{8*pi},100000);

\node[
    anchor=north west
]
at (axis cs:{8*pi*1.03},50000)
{
    $\displaystyle d=4\lambda$
};

% Feldbereiche
\node[
    anchor=south,
    font=\scriptsize\bfseries
]
at (axis cs:0.1,130000)
{Reaktives Nahfeld};

\node[
    anchor=south,
    font=\scriptsize\bfseries
]
at (axis cs:5,130000)
{Strahlendes Nahfeld};

\node[
    anchor=south,
    font=\scriptsize\bfseries
]
at (axis cs:50,130000)
{Fernfeld};

% Erläuterung des elektrischen Dipols
\node[
    anchor=west,
    fill=white,
    draw=DARCorange,
    align=left
]
at (axis cs:0.025,40000)
{
    Elektrischer Dipol
};

% Erläuterung der magnetischen Loop
\node[
    anchor=west,
    fill=white,
    draw=DARCgreen,
    align=left
]
at (axis cs:0.025,3)
{
    Magnetische Loop
};

\end{axis}

\end{tikzpicture}